\documentclass[11pt,reqno]{amsart}
\usepackage[T1]{fontenc}\usepackage{lmodern}
\usepackage{amsmath,amssymb,amsthm,mathtools}
\usepackage[margin=1in]{geometry}\usepackage{microtype}\usepackage{enumitem}
\usepackage[colorlinks=true,linkcolor=blue!55!black]{hyperref}
\theoremstyle{plain}\newtheorem{theorem}{Theorem}\newtheorem{proposition}[theorem]{Proposition}
\newtheorem{lemma}[theorem]{Lemma}\newtheorem{corollary}[theorem]{Corollary}
\theoremstyle{definition}\newtheorem{definition}[theorem]{Definition}
\theoremstyle{remark}\newtheorem{remark}[theorem]{Remark}
\newcommand{\KR}{K_R}\newcommand{\KF}{K_F}\newcommand{\KM}{K_M}
\newcommand{\KRo}{K_R^{\circ}}\newcommand{\KFo}{K_F^{\circ}}
\newcommand{\CM}{C_M}\newcommand{\CF}{C_F}\newcommand{\CR}{C_R}\newcommand{\Dbar}{\bar D}
\newcommand{\bp}{\beta+p}\newcommand{\E}{\mathbb{E}}\newcommand{\phio}{\varphi^{\circ}}
\newcommand{\qo}{q^{\circ}}\newcommand{\FKR}{F_{\!\KR}}\newcommand{\eqdef}{\vcentcolon=}
\setlength{\parskip}{3pt}
\begin{document}

\title[Prices versus Quantities for Greening Flexible Computing]
{Prices versus Quantities for Greening Flexible Computing:\\
Allocation Equivalence, Revenue Divergence, and the Long-Run Ranking}
\author{Team A1 --- Role 1}
\date{Draft for \S5.3}

\begin{abstract}
Two instruments can stop a flexible computing load from burning fossil fuel: a carbon tax
$\tau$ on fossil output, and a dispatch cap that permits the load to run only on realized
renewable surplus. At their respective greening thresholds the two are
\emph{allocation-equivalent}: identical computing, identical zero carbon from the load, and
--- as the Guide observes --- an identical marginal-entrant condition. We show the
equivalence is fragile in exactly the dimension the policy debate cares about. The tax
delivers strictly positive revenue to renewables in the hours the load induces; the cap
delivers none. Because that revenue enters the renewable first-order condition additively,
the two instruments induce \emph{different} equilibrium renewable capacity, and the
resulting consumer carbon is strictly ordered. Short-run equivalence therefore does not
survive endogenous investment, and the ranking favours the price instrument for a reason
Weitzman's calculus does not see: not relative slope uncertainty, but who is paid. We also
close the boundary problem: at the greening tax exactly, producers are indifferent and the
equilibrium quantity is a correspondence, not a point; the $\tau\downarrow\bar\tau$ limit
selects the fossil-off equilibrium, which is the selection the comparison requires.
\end{abstract}
\maketitle

\section{Setting}
We work in the deregulated (spot) environment, where the relevant computing margin is
$R-c-\tau$ rather than $R-p-\tau$. A tax $\tau$ per fossil MWh raises the effective fuel
cost to $c+\tau$; every formula of the baseline applies under that substitution. In
particular the no-load fossil FOC $(\bp-c-\tau)[1-F(\KFo)]=\CF$ gives the benchmark
reliability
\[
  \phio(\tau)=1-\frac{\CF}{\bp-c-\tau},
\]
and we maintain $\bp>R+\CF$ so this stays positive up to
\[
  \bar\tau\eqdef R-c,
\]
the tax at which a miner-valued MWh no longer covers taxed fuel. Write $\qo(\tau)=F(0;\KRo(\tau))$
for the no-load surplus probability. The \emph{cap} (renewable-only contract) is
$d(x)\le(-x)_+$: the load may run only on realized surplus.

\section{Allocation equivalence at the greening thresholds}

\begin{lemma}[Viability splits into a taxed and a tax-exempt part]\label{lem:split}
Under the tax the first entrant is dispatched whenever spare capacity exists. In surplus
states (probability $\qo$) it buys at $P=0$ and keeps $R$; in fossil-topped states
(probability $\phio(\tau)-\qo$) it buys at $c+\tau$ and keeps $R-c-\tau$. Hence
\[
  \bar\CM(\tau)=\underbrace{(R-c-\tau)\big[\phio(\tau)-\qo\big]}_{\text{fossil-hour part}}
  \;+\;\underbrace{R\,\qo}_{\text{surplus part}} .
\]
\end{lemma}

\begin{theorem}[Allocation equivalence]\label{thm:equiv}
At $\tau=\bar\tau$ the fossil-hour part of $\bar\CM$ vanishes and
$\bar\CM(\bar\tau)=R\,\qo$, which is exactly the cap's entry threshold: under either
instrument the first entrant earns $R$ in surplus hours only. Consequently, at
$(\tau=\bar\tau)$ versus (cap), holding capacities fixed:
\textup{(i)} the computing dispatch schedule coincides ($d=(-x)_+\wedge\KM$),
\textup{(ii)} carbon attributable to the load is zero under both, and
\textup{(iii)} the marginal-entrant condition is the same, so the same computing capacity
enters. The two instruments admit the same computing at the margin.
\end{theorem}
\begin{proof}
(i) At $\tau=\bar\tau$ fossil-served computing earns $R-c-\bar\tau=0$ and is weakly
dominated; the selection of Theorem~\ref{thm:selection} removes it, leaving surplus-only
dispatch, which is the cap's schedule. (ii) Immediate from (i): the load draws no fossil.
(iii) Lemma~\ref{lem:split} at $\tau=\bar\tau$ gives $\bar\CM(\bar\tau)=R\qo$; under the cap
the entrant runs only in surplus hours and earns $R$ there, so its threshold is $R\qo$ as well.
\end{proof}

\section{The boundary: indifference, and the selection that resolves it}

At $\tau=\bar\tau$ producers are exactly indifferent about serving computing from fossil, so
the market-clearing \emph{quantity} is not pinned by optimization alone. This is the gap
flagged in review, and it must be closed before Theorem~\ref{thm:equiv} can be used.

\begin{theorem}[Price is unique, quantity is a correspondence, and the limit selects]\label{thm:selection}
Fix capacities and a state $x\in(-\KM,0]$ with fossil headroom.
\begin{enumerate}[label=\textup{(\alph*)},leftmargin=2.2em]
  \item For $\tau<\bar\tau$: fossil tops up at marginal cost $c+\tau<R$, so $P=c+\tau$ and
        fossil-served computing is strictly profitable; the equilibrium is unique.
  \item At $\tau=\bar\tau$: $c+\tau=R$, so the clearing price is \emph{unique}, $P=R$; but
        every fossil-to-computing quantity in $[0,\,\KM-(-x)_+]$ is consistent with producer
        indifference and market clearing. The equilibrium \emph{quantity} is a correspondence.
  \item For $\tau>\bar\tau$: fossil-served computing is strictly unprofitable, so the unique
        equilibrium has fossil-to-computing quantity $0$.
\end{enumerate}
Consequently the correspondence at $\bar\tau$ is the upper limit of (a) and the lower limit
of (c), and the refinement $\tau\downarrow\bar\tau$ selects the fossil-off point. We adopt
that selection throughout; equivalently, one may state all results for $\tau>\bar\tau$ and
pass to the limit, which is the conservative reading.
\end{theorem}
\begin{proof}
(a) and (c) are immediate from the sign of $R-c-\tau$. For (b), at $c+\tau=R$ a fossil
generator earns exactly zero on a computing MWh and a miner earns exactly zero, so any
split clears; the price cannot exceed $R$ (demand vanishes above it) nor fall below $c+\tau$
(supply vanishes below it), and these coincide, pinning $P=R$. The limit statement follows
because the fossil-off quantity is the unique selection on $(\bar\tau,\infty)$ and the
correspondence is upper hemicontinuous at $\bar\tau$.
\end{proof}

\begin{remark}
The same argument disposes of the companion objection that ``the price never leaves $0$''
under the cap. Under the cap, miner demand in surplus states equals the zero-marginal-cost
surplus exactly, so any $P\in[0,R]$ clears; the competitive selection is $P=0$ because
supply is perfectly elastic at zero up to the surplus. That is a selection, not an
implication, and we flag it as such.
\end{remark}

\section{Revenue divergence, and why the equivalence dies in the long run}

\begin{proposition}[Revenue divergence]\label{prop:revenue}
Per unit of renewable capacity per hour, the two instruments deliver
\[
  \text{tax:}\quad \E\big[\tilde S\,\min\{c+\tau,R\}\,\mathbf 1_{\{-\KM<\tilde x\le 0\}}\big]>0,
  \qquad\text{cap:}\quad 0 .
\]
The tax's pass-through to renewable revenue is nondecreasing in $\tau$ and saturates at
$\bar\tau$; the cap leaves surplus-hour prices at zero by construction, because miner demand
there exactly equals the zero-marginal-cost surplus.
\end{proposition}

This is the crux. The two instruments buy the same allocation at different prices, and the
difference is paid to renewable capacity.

\begin{theorem}[Long-run non-equivalence and a strict ranking]\label{thm:longrun}
Let $\KR^{\mathrm{tax}}$ and $\KR^{\mathrm{cap}}$ denote equilibrium renewable capacity under
the greening tax and under the cap. The renewable first-order conditions differ by exactly
the strictly positive term of Proposition~\ref{prop:revenue}; hence, whenever the renewable
marginal value is strictly decreasing \textup{(}the standing single-peakedness\textup{)},
\[
  \boxed{\ \KR^{\mathrm{tax}}>\KR^{\mathrm{cap}},\qquad\text{and therefore}\qquad
  \E[\text{fossil}^{\mathrm{tax}}]<\E[\text{fossil}^{\mathrm{cap}}].\ }
\]
Allocation equivalence at fixed capacities does \emph{not} survive endogenous investment:
the instruments are strictly ranked, and the price instrument dominates.
\end{theorem}
\begin{proof}
Both instruments induce the same computing dispatch (Theorem~\ref{thm:equiv}), so the
consumer-side terms of the renewable FOC coincide. The tax adds the surplus-hour revenue of
Proposition~\ref{prop:revenue}, which is strictly positive whenever $\KM>0$ and
$\Pr(-\KM<\tilde x\le0)>0$. Evaluating the tax economy's renewable marginal value at
$\KR^{\mathrm{cap}}$ therefore exceeds $\CR$; by single-peakedness the tax optimum lies
strictly to the right. The carbon conclusion follows from
$e(\KR)=\int_0^{\KF(\KR)}(1-F)$ being strictly decreasing in $\KR$.
\end{proof}

\begin{corollary}[What the ranking is \emph{not}]
The ranking is not an efficiency comparison in Weitzman's sense: it does not turn on the
relative slopes of marginal benefit and marginal cost under uncertainty. Both instruments
implement the same allocation. The wedge is \emph{distributional} --- who receives the
scarcity rent in surplus hours --- and it becomes an efficiency difference only because that
rent is also the marginal return to renewable investment. A quantity instrument that
rebated the same rent to renewables would replicate the tax.
\end{corollary}

\begin{remark}[Policy reading]
New York's moratorium on fossil-fired proof-of-work mining, which exempts renewable-powered
operations, is a cap. Theorem~\ref{thm:longrun} says that such a rule buys the same
short-run decarbonization as the equivalent carbon price and forgoes the renewable
build-out the price would have financed. The cap is the feasible near-term instrument at a
known cost, and the cost is quantified here rather than asserted.
\end{remark}

\end{document}
